\section{本章小结}
\label{sec:lop-summary}

本章把函数到函数的变换写成线性算子，并始终把定义域、取值空间和两侧的范数当作算子的一部分。恒等、投影与乘以有界函数自然定义在整个空间上；积分公式要由核、测度与输入空间共同保证有意义；求导只能定义在可导函数组成的真子空间上。说一个算子“有界”，就是说存在统一常数$C$使$\lVert Tx\rVert\leq C\lVert x\rVert$；对线性算子，这等价于连续与Lipschitz连续，最小的这种常数就是算子范数，它统一控制输入差造成的输出差。

同一个公式在不同范数下可以有完全不同的行为。求导在一致范数下无界，在$C^1$范数下有界；平方可积核给出有界积分算子，且$\lVert T\rVert\leq\lVert k\rVert_{L^2}$。包含映射把“同一个函数在不同空间中的大小”变成算子的有界性：有限测度上$L^q$连续嵌入$L^p$（$p\leq q$），计数测度上方向相反；改变测度时先由绝对连续性保证等价类的对应，再由密度的上下界决定有界性与范数等价。算子列的范数收敛给出统一误差，强算子收敛只给出逐输入误差，截断投影说明两者确实不同。

Hilbert空间中，伴随把输出方向拉回输入空间，给出$(\operatorname{Ran}T)^\perp=\ker T^*$与$\overline{\operatorname{Ran}T}=(\ker T^*)^\perp$；像未必闭，与$\ker T^*$正交的输出只能被逼近，未必真正可达。自伴是同一空间上$T=T^*$的附加结构，正性还要求二次型非负，二者都不能由核在个别点上的符号判断。紧算子把有界输入列变成有收敛子列的输出列，把弱收敛改善为范数收敛；在Hilbert空间中它们恰好是有限秩算子的算子范数极限，平方可积核的积分算子都是紧算子，核的$L^2$误差控制算子范数误差。

完备性把较弱的信息升级为统一控制：一致有界原理由逐输入有界得到统一的范数上界，开映射定理与有界逆定理由满射与双射得到稳定的逆，闭图定理由图的闭性得到有界性。这些结论都需要定义域与取值空间完备，定义在真子空间上的算子须先换用使定义域完备的图范数。第~\ref{chap:operator-spectra}章在这些基础上研究$T-\lambda I$何时可逆，并为紧自伴算子与一般紧算子建立特征展开和奇异分解；给定输出求解算子方程的可解性与稳定性，留到第~\ref{chap:linear-operator-equations}章。
